\section{Metodología de Respuesta a Fallas}

Para garantizar una respuesta rápida y efectiva a las fallas de la red, se seguirá la siguiente metodología basada en el ciclo de vida de la gestión de incidentes:

\begin{enumerate}
    \item \textbf{Detección y Registro:} Las fallas pueden ser detectadas por el sistema de monitoreo proactivo o reportadas por los usuarios. En ambos casos, se registrará un ticket en el sistema de soporte con toda la información relevante.
    \item \textbf{Clasificación y Priorización:} Cada incidente se clasificará según su impacto (número de usuarios afectados) y urgencia (criticidad de los servicios afectados). Esto determinará la prioridad para su atención.
    \begin{itemize}
        \item \textbf{Prioridad Alta:} Falla general de la red, afectando a un gran número de usuarios o servicios críticos.
        \item \textbf{Prioridad Media:} Falla parcial que afecta a un grupo de usuarios o a un servicio no crítico.
        \item \textbf{Prioridad Baja:} Problema de conectividad de un solo usuario o degradación menor del servicio.
    \end{itemize}
    \item \textbf{Diagnóstico y Solución:} El equipo de soporte técnico investigará la causa raíz del problema. Una vez identificada, se aplicará la solución correspondiente, que puede ir desde un reinicio de equipo hasta una reconfiguración de la red.
    \item \textbf{Recuperación y Cierre:} Después de aplicar la solución, se verificará que el servicio se haya restablecido por completo y que el problema esté resuelto. Se documentará la solución en el ticket y se procederá a su cierre, notificando al usuario.
    \item \textbf{Análisis Post-Incidente:} Para incidentes de alta prioridad, se realizará un análisis posterior para identificar la causa fundamental y definir acciones preventivas que eviten su recurrencia en el futuro.
\end{enumerate}
